\iftoggle{paper}
{


This document presents the thesis topic proposal for obtaining the Master's degree in Computer Science: Emphasis in Computer Science at the Tecnológico de Costa Rica.
The proposed topic is centered in the idea of the potential use of machine learning techniques to analyze the patterns in acoustic echoes produced inside small rooms to estimate the location of certain objects in a tridimensional space. Two neural networks architectures are proposed for this task, a Convolutional Neural Networks that despite being an architecture mainly for images is commonly used in literature for acoustic inference, and an Echo State Machine as a new architecture with potential to address directly the acoustic signals. The effect of beamforming is also considered as one of the elements to investigate as a solution to the multi-path acoustic superposition. The design for the research is presented, the required experimental setup is discussed as well as the methodology to measure, process and evaluate the results, all this with the objective to evaluate the potential use of sound echoes in the audible frequency as a mean to locate object in an enclose room.        
 
}%

